\chapter{Chemical reaction networks as programs}

Place four molecules of species $A$ and four of $B$ in an ideal well-mixed
compartment. Permit only $A+B\rightarrow W$. Predict the final counts.
Now predict the time at which the last pair disappears. The first question
has an exact answer without a rate constant; the second does not.
This distinction is the starting point for treating chemistry as a program.

Chapter 19 reconstructed one strand-exchange mechanism. Here we construct
the language in which many such mechanisms can be composed. We will specify
what a reaction consumes, what it produces, when it can occur, how often it
occurs, and what counts as a correct output. A formal reaction is not yet
a realizable DNA design. Our examples are mathematical references, with
assigned rates and no new experimental measurements.

\section{Name the state before naming the algorithm}

A \Term{species} is a distinguishable molecular type at the chosen level of
description. It may be a free strand, a duplex, a bound complex, or a purely
formal signal awaiting physical implementation. Fix an ordered list
$S_1,\ldots,S_s$. A discrete state is a vector
$n\in\mathbb N^s$ of molecule counts. A continuous state is a vector
$x\in\mathbb R_{\geq0}^s$ of concentrations. Their entries are not interchangeable.

For $r$ reactions, let $R,P\in\mathbb N^{s\times r}$ contain reactant and product
stoichiometries. Column $j$ represents
\begin{equation}
 \sum_i R_{ij}S_i\longrightarrow\sum_i P_{ij}S_i,\qquad
 N=P-R .
 \label{eq:stoich}
\end{equation}
$N$ is the \Term{stoichiometric matrix}. It describes the net state change,
but not the enabling condition. Reaction $j$ is enabled at $n$ precisely when
$n_i\geq R_{ij}$ for every $i$. Its firing then produces $n+N_{\cdot j}$.
The reaction $A+F\rightarrow A+Y$ has zero net change in $A$, but still
requires $A$. Deleting the catalyst from $R$ because its row of $N$ is zero
changes the program.

For $A+B\rightarrow W$, in species order $(A,B,W)$,
$R=(1,1,0)^\top$, $P=(0,0,1)^\top$ and $N=(-1,-1,1)^\top$.
For $2A\rightarrow W$, the reactant column is instead $(2,0,0)^\top$.
Two molecules of one species are not the same state requirement as one
molecule each of two species.

\Plate{01-coordinates}{Columns define reaction events, not a chronological
instruction list. The two reactions have different enabling conditions even
though both consume two molecules and create one complex. Labels identify
species independently of color.}{fig:coordinates}

\Listing{network}
The constructor takes ownership of copies and marks them read-only. It limits
stoichiometric coefficients to 100 to keep this small numerical reference
within its intended scope. A frozen Python container alone would not make
the arrays immutable. The companion code also validates state shapes, finite
values and integer counts; counts above $10^9$ are outside its declared domain.
These are implementation limits, not limits of CRN theory.

\section{Conservation is a proof obligation, not a visual impression}

Let $c\in\mathbb R^s$ satisfy $c^\top N=0$. Then every firing preserves
$c^\top n$, because
\begin{equation}
 c^\top(n+N_{\cdot j})=c^\top n+c^\top N_{\cdot j}=c^\top n.
\end{equation}
Such a vector is a \Term{linear invariant}. It can express material inventory
or a logical difference. For annihilation, both $A+W$ and $B+W$ are invariant,
and subtracting them gives the invariant $A-B$. The latter can be negative;
it is not itself a nonnegative material inventory.

Return to the strand-displacement reaction $I+G\rightarrow W+O$.
In order $(I,G,W,O)$, the three strand inventories form the matrix
\begin{equation}
 C=
 \begin{pmatrix}
 1&0&1&0\\
 0&1&1&0\\
 0&1&0&1
 \end{pmatrix},
 \qquad
 N=\begin{pmatrix}-1\\-1\\1\\1\end{pmatrix},
 \qquad CN=0.
\end{equation}
The rows count invader, substrate and incumbent strands. In particular,
$G$ and $W$ each contain two strands. Counting species equally happens to
give two reactants and two products here, but that coincidence is not the
reason the molecular mechanism conserves material.

\Plate{02-inventory}{The formal exchange column is backed by continuous
covalent backbones. Arrowheads on strands mark the $3'$ direction; short
vertical marks denote pairing. The reaction arrow changes occupancy, not
strand identity. The matrix counts the three physical strand types.}{fig:inventory}

In a closed physical system, atom inventories must also balance. A formal
signal species need not represent a single molecule, however: a compiler may
use auxiliary gates and fuels that are absent from the signal-level equations.
An apparent source or sink in the signal model therefore requires an explicit
reservoir in its physical interpretation. A left-nullspace calculation is a
useful audit, not a complete chemical realizability theorem.

\section{Separate the jump rule from the kinetic law}

Under deterministic mass action, reaction $j$ has flux
\begin{equation}
 v_j(x)=k_j\prod_i x_i^{R_{ij}},\qquad
 \frac{dx}{dt}=Nv(x).
 \label{eq:ode}
\end{equation}
If concentrations have unit $U$ and time has unit seconds, a reaction of
total order $q_j=\sum_iR_{ij}$ has $k_j$ in
$U^{1-q_j}\mathrm{s}^{-1}$. Flux has unit $U\,\mathrm{s}^{-1}$.
This convention gives flux $kx_A^2$ for $2A\rightarrow W$ and
$\dot x_A=-2kx_A^2$; the factor two is in stoichiometry.

At a boundary where $x_i=0$, every mass-action term consuming $S_i$ is zero,
while terms producing it are nonnegative. Thus the vector field cannot
point out of the nonnegative orthant. This is a property of the differential
equation, not of every numerical integrator. A large explicit Euler step can
still yield negative values. Likewise,
$d(c^\top x)/dt=c^\top Nv(x)=0$ proves continuous conservation.

For molecule counts, let $\Omega$ convert concentrations to counts:
$n_i=\Omega x_i$. For molar concentrations,
$\Omega=N_{\!A}V$ when $V$ is expressed in litres. The same macroscopic
$k_j$ convention gives the propensity
\begin{equation}
 a_j(n)=k_j\Omega^{1-q_j}
 \prod_i(n_i)_{R_{ij}},
 \qquad
 (n)_r=n(n-1)\cdots(n-r+1).
 \label{eq:propensity}
\end{equation}
Set $(n)_0=1$ and $(n)_r=0$ if $n<r$. Propensity is an event hazard in
$\mathrm{s}^{-1}$, not a concentration flux. In particular,
$a=k n_A n_B/\Omega$ for $A+B\rightarrow W$ and
$a=k n_A(n_A-1)/\Omega$ for $2A\rightarrow W$.

Some stochastic conventions write a rate constant times $\binom{n_A}{2}$.
That is compatible only after converting the rate constant: here it would
be $2k/\Omega$. Reusing the same numeric coefficient with both formulas
introduces a factor-of-two error. With only one $A$, the correct propensity
is zero, whereas replacing the falling factorial by $n_A^2$ permits an
impossible event.

\Listing{propensities}
The empty product for a zero-order source is one, giving hazard $k\Omega$.
Such a source is useful in an open-system abstraction; it is not spontaneous
creation of physical atoms. The drift function in the companion reference
implements Equation~\ref{eq:ode}. The firing function below implements only
stoichiometric enabledness: a zero kinetic rate can still label a formally
enabled reaction, but a stochastic scheduler must never select it.
\Listing{fire}

\section{Prove a small reaction program}

Consider the one-reaction network $A+B\rightarrow W$ with initial
counts $(a,b,0)$ and strictly positive rate. We use residual $A$ and $B$
as a dual-rail representation of a signed difference. The following result
does not require a stochastic rate model.

\begin{proposition}[Annihilation computes a residual difference]
Every maximal firing sequence terminates after $\min(a,b)$ firings at
\[
 (\max(a-b,0),\ \max(b-a,0),\ \min(a,b)).
\]
\end{proposition}
\begin{proof}
The nonnegative integer ranking function $n_A+n_B$ decreases by two at
every firing, so no firing sequence is infinite. A terminal state must have
$n_A=0$ or $n_B=0$, since otherwise the sole reaction is enabled.
The invariants $n_A+n_W=a$ and $n_B+n_W=b$ then force the displayed state.
\end{proof}

This is a terminal-output theorem. It does not tell an observer when the
state has become terminal. A low signal might mean no molecules, slow
reaction, failed reporting, or a concentration below detection. Equal inputs
produce no residual signal on either rail; a practical equality decision
requires an observation policy, not just a proof of arithmetic.

For example, $(8,4,0)$ terminates at $(4,0,4)$. Interpreting $W$ as the
answer instead computes $\min(a,b)$, not $a-b$. The same reaction system
supports different functions under different input and output encodings.
Specify the decoder as part of the program.

\section{Correct completion can still have an expensive tail}

In the stochastic model the only possible event has hazard
$k(a-j)(b-j)/\Omega$ after $j$ firings. Each holding time is exponential;
successive holding times are independent by the Markov property. Consequently
\begin{align}
 \mathbb E[T]&=\frac{\Omega}{k}
 \sum_{j=0}^{\min(a,b)-1}\frac{1}{(a-j)(b-j)},\\
 \operatorname{Var}(T)&=\frac{\Omega^2}{k^2}
 \sum_{j=0}^{\min(a,b)-1}\frac{1}{(a-j)^2(b-j)^2}.
 \label{eq:timing}
\end{align}
If either input is zero, the empty sum gives zero time. For equal inputs
$a=b=m$, the mean is $(\Omega/k)\sum_{\ell=1}^m1/\ell^2$.
The final pair alone contributes $\Omega/k$, even though the first event
has mean $\Omega/(km^2)$. Fast initial progress does not certify a short tail.

\Listing{completion_moments}
\begin{center}\input{\Generated/completion.tex}\end{center}
The table assigns $k=1$ and $\Omega=1$ in consistent reference units. It is
not a laboratory timing estimate. Holding volume fixed while increasing $m$
also increases concentration. Under a fixed initial concentration
$c_0=m/\Omega$, instead $\Omega=m/c_0$ and the mean grows linearly with $m$,
up to the bounded factor $\sum_{\ell=1}^m1/\ell^2$. These are different
scaling experiments.

\Plate{03-tail}{An assigned-rate stochastic path reaches zero after four
events, while the equal-concentration ODE tail remains positive at every
finite time. The count-equivalent ODE is $4/(1+4t)$ for $k=\Omega=1$.
The staircase is one seeded sample, not an ensemble mean or measured trace.}{fig:tail}

For equal initial concentrations, the ODE gives
$x_A(t)=c_0/(1+kc_0t)$. It never reaches exactly zero at finite time.
Moreover, $d\mathbb E[n_A]/dt=-(k/\Omega)\mathbb E[n_An_B]$;
the product expectation is generally not the product of the expectations.
Replacing it by $\mathbb E[n_A]\mathbb E[n_B]$ is a closure approximation.
Large-system limits can justify deterministic trajectories under appropriate
conditions, but do not make every finite-copy mean exactly obey mass action.

A conservative deadline bound follows from Markov's inequality:
$\Pr(T\geq t)\leq\mathbb E[T]/t$ for $t>0$. This is often loose. It is still
more honest than interpreting a mean as a guaranteed maximum. Our tests
compare the timing formulas with an independent backward recursion and
check a seeded ensemble against a declared five-standard-error tolerance.
Those tests validate the implementation of the stipulated process.

\section{Compose two modules and watch the shared resource}

Suppose two abstract catalytic modules share one fuel species:
\begin{equation}
 A+F\xrightarrow{k_1}A+Y,\qquad
 B+F\xrightarrow{k_2}B+Z.
 \label{eq:fuel}
\end{equation}
Here catalysts are unchanged by either reaction. With fixed catalyst
concentrations, write $\alpha=k_1[A]$ and $\beta=k_2[B]$, each in inverse
seconds. Let initially $F=F_0$ and $Y=Z=0$. The material-accounting invariant
is $F+Y+Z=F_0$, and
\begin{align}
 F(t)&=F_0e^{-(\alpha+\beta)t},\\
 Y(t)&=F_0\frac{\alpha}{\alpha+\beta}
             (1-e^{-(\alpha+\beta)t}),\\
 Z(t)&=F_0\frac{\beta}{\alpha+\beta}
             (1-e^{-(\alpha+\beta)t}).
 \label{eq:fuel-sol}
\end{align}
For $\alpha+\beta=0$, no conversion occurs. The formulas follow by first
solving $\dot F=-(\alpha+\beta)F$ and then integrating $\dot Y=\alpha F$
and $\dot Z=\beta F$.

\Plate{04-fuel}{Two individually valid catalytic modules compete when they
share a reservoir. Dashed return arcs identify catalysts; solid conversion
arrows consume fuel. The equations are a formal catalytic model, not a claim
that one strand-displacement gate realizes either entire cycle.}{fig:fuel}

\Listing{shared_fuel}
\begin{center}\input{\Generated/fuel.tex}\end{center}
For $F_0=100$, $\alpha=0.1$ and $\beta=0.3$, the limiting products are
25 and 75. Module one in isolation would convert all 100 units to $Y$;
composing it with module two changes that answer even if their signals
never cross-hybridize. This is resource coupling, not sequence crosstalk.
Duplicating $F_0$ in two separate simulators silently duplicates material.

In a finite-copy version with constant catalysts, the final allocation of
$F_0$ integer fuel molecules to $Y$ is binomial with parameter
$\alpha/(\alpha+\beta)$. The mean split agrees with the formula, but a single
compartment need not hit that mean exactly. A physical realization also
needs compatible molecular compositions or hidden waste species. The
formal invariant names a fuel-equivalent inventory, not every atom.

\section{What must a molecular compiler preserve?}

Let $z$ denote a detailed state including intermediates and auxiliaries,
and let $x=Lz$ be a declared linear coarse representation. Sometimes $L$
counts both free and bound forms of a signal; sometimes only free signal
is the interface quantity. These choices have different kinetics.
If $\dot z=f(z)$, an exact closed coarse model $\dot x=g(x)$ requires
\begin{equation}
 Lf(z)=g(Lz)
 \label{eq:closure}
\end{equation}
throughout the relevant state set. Thus two detailed states with the same
$Lz$ must induce the same coarse derivative. A hidden bound complex can
violate this condition even when both states have equal free-signal counts.

\Plate{05-projection}{A projection merges detailed states only with an
explicit semantic contract. Equal projected signal does not force equal
release flux when intermediate occupancy differs. Approximate closure needs
a domain and error criterion; dropping intermediate names is not a proof.}{fig:projection}

Most useful reductions are approximate rather than exact. One may require,
for declared initial conditions and observation interval $[0,T]$,
$\sup_{0\leq t\leq T}\|Lz(t)-x_{\rm target}(t)\|\leq\varepsilon$.
Specify the norm, concentration scale, admissible inputs, auxiliary inventories
and rate-separation assumptions. A good endpoint match alone need not satisfy
a trajectory-level contract.

Soloveichik, Seelig and Winfree construct strand-displacement cascades that
approximate coupled formal reaction dynamics \cite{soloveichik}. That result
motivates the compiler interface; it does not eliminate auxiliary species,
finite fuel or the need to check the regime of approximation.
Our four-line mass-action evaluator is not an implementation of that compiler.

\section{Stable computation and current research boundaries}

Our annihilation proof is stronger than a favorable simulated run: every
maximal execution of its one reaction has the stated terminal result.
For general networks, stable output need not mean the entire network has
halted. The output may stop changing while other reactions continue.
Reachability, eventual stabilization, deadlines and noise tolerance are
separate obligations.

Chen, Doty and Soloveichik characterize deterministic count-valued function
computation in their discrete stable-computation model by semilinear graphs
\cite{chen}. This is a model-specific characterization, not a claim that all
continuous chemical dynamics are semilinear or that no probabilistic CRN can
simulate more general computation. It also does not grant a physical device
unbounded molecule counts or perfect detection.

A 2026 result by Kini and Doty studies a different boundary: reactions may
reverse before a cutoff, but only forward reactions remain afterwards
\cite{kini}. Their reverse-robust model retains semilinear predicate and
function computability. Importantly, that cutoff assumption is not persistent
thermal reversibility at equilibrium. For our annihilation example, adding
$W\rightarrow A+B$ forever preserves $A-B$ but destroys terminal zero on one
rail. Invariance by itself therefore does not prove stable readout.

For an AI-assisted molecular-design pipeline, the CRN is an intermediate
specification with testable obligations. Store the reaction convention and
units; reject disabled-event bugs; audit inventories; compare detailed and
coarse trajectories; then validate experimentally with appropriate controls.
A generated reaction list is neither a proof of implementability nor evidence
of useful biological behavior. Chapter 21 uses these obligations to examine
thresholds, restoration and fan-out in composed circuits.

\section{Exercises with worked answers}
\begin{enumerate}
\item Write $R,P,N$ for $2A\rightarrow W$ in order $(A,W)$.
\textbf{Answer:} $(2,0)^\top$, $(0,1)^\top$, $(-2,1)^\top$.
The inventory $A+2W$ is conserved.
\item Why can $N$ not determine whether a catalyst is required?
\textbf{Answer:} $F\rightarrow Y$ and $A+F\rightarrow A+Y$ have the same
net-change column on $(A,F,Y)$ but different reactant columns.
\item Fire annihilation from $(7,3,0)$ to termination.
\textbf{Answer:} Three firings give $(4,0,3)$. Residual rails encode $+4$;
$W$ encodes the minimum, 3.
\item Compute the $2A\rightarrow W$ hazard at $n_A=3$, $k=2$, $\Omega=10$.
\textbf{Answer:} $2(3)(2)/10=1.2\ {\rm s}^{-1}$ under our convention.
The binomial coefficient version requires coefficient $0.4\ {\rm s}^{-1}$.
\item Show that the continuous material invariant is exact.
\textbf{Answer:} Multiply $\dot x=Nv(x)$ by a left-null vector;
$c^\top\dot x=0$. This says nothing about numerical time-step error.
\item Find mean and variance of completion from $a=b=2$, $k=\Omega=1$.
\textbf{Answer:} Mean $1/4+1=1.25$; variance $1/16+1=1.0625$.
\item At fixed initial concentration, why does the last-pair mean grow with
input size? \textbf{Answer:} $\Omega=m/c_0$ grows, so its hazard $k/\Omega$
shrinks. Volume is part of the cost model.
\item Give a deadline guaranteeing the Markov upper bound is at most $0.05$
when $\mathbb E[T]=2$ seconds.
\textbf{Answer:} $t=40$ seconds; this is a bound, not a measured quantile.
\item Compute the limiting shared-fuel split for $\alpha=\beta>0$.
\textbf{Answer:} Half to each output in the deterministic model; finite-copy
allocation fluctuates around that split.
\item Does preserving $A-B$ prove stable output after adding reverse reactions?
\textbf{Answer:} No. Reverse events recreate both rails while preserving
their difference; a zero-rail output predicate can cease to hold.
\item What falsifies exact projection closure?
\textbf{Answer:} Find $z_1,z_2$ with $Lz_1=Lz_2$ but $Lf(z_1)\ne Lf(z_2)$.
Then no single-valued $g$ satisfies Equation~\ref{eq:closure} on both.
\item Design a composition validation study.
\textbf{Rubric:} State signal encodings, units and output criterion; compare
isolated modules with the full shared-inventory model; include no-input and
resource-limited cases; audit strand and fuel inventories; specify trajectory
error and observation window; distinguish simulation from calibrated evidence.
\end{enumerate}

\section{Vocabulary and the next interface}
\Term{Propensity}: instantaneous event hazard for a discrete count state.
\Term{Stable output}: an output that cannot subsequently change under the
model's allowed continuation.
\Term{Projection closure}: the coarse state determines its own evolution
without additional hidden state.
\Term{Resource coupling}: dependence introduced by a shared finite inventory.
The next chapter asks how a circuit restores distinguishable signals while
paying those material and timing costs.
